\section{High-weight cosets and code extensions}
\label{sec:high-weight-cosets}

We now prove Theorem~\ref{thm:main} and record its two coding consequences.
The reader needs only the conclusion of the preceding sections: outside
\(\cP_r\cup\cM^{\max}_{r,p}\), there is a degree-\((r-2)\) split squarefree
locator supported on \(S\).

\subsection{Coset weights}

Let the columns of \(H_S\) be \(w_a\nu_{r-1}(a)\), \(a\in S\), with every
\(w_a\ne0\).  Rescaling an error value absorbs \(w_a\), so
\begin{equation}\label{eq:grs-coset-weight}
 d_S(f)=\min\bigl\{|T|:T\subseteq S,\quad
 f\in\langle\nu_{r-1}(T)\rangle\bigr\}.
\end{equation}
Thus arbitrary GRS multipliers change neither coset weight nor the projective
syndrome geometry.  The Hankel criterion identifies a degree-\((r-2)\) split
squarefree member of \(W_f\), supported on \(S\), with a representation in
\eqref{eq:grs-coset-weight} of size \(r-2\).

\begin{proof}[Proof of Theorem~\ref{thm:main}]
The pointed simultaneous-marker theorem gives
\(d_S(f)\leq r-2\) outside
\(\cP_r\cup\cM^{\max}_{r,p}\), with the two stated thresholds.  If
\(p>r-1\), no two adjacent entries of Pascal row \(r-2\) vanish, so the Lucas
carrier is empty.

The rational points of \(\cP_r\) have four disjoint strata: curve points,
interiors of rational split secants, off-curve tangent points, and rational
points of conjugate secants.  When \(s=0\), the Seroussi--Roth--D\"ur radius
gate gives radius \(r-1\).  Curve points and split-secant interiors have
weights one and two, while the full-support rank-two calculation
\cite[Thms.~I.5--I.6]{ZWK2020} gives weight \(r-1\) on every off-curve tangent
point and rational point of a conjugate secant.  These are therefore exactly
the maximum-weight strata.

Suppose \(s>0\).  If \(a\in A\), the arc property says that
\(\nu_{r-1}(a)\) is outside the span of every \(r-1\) retained columns, while
any \(r\) retained columns span the syndrome space.  Hence
\(d_S(\nu_{r-1}(a))=r\).  Seroussi and Roth's unique-extension theorem
\cite[Thm.~1]{SeroussiRoth1986} applies when
\[
 2\leq r\leq |S|-\left\lfloor\frac{q-1}{2}\right\rfloor.
\]
Indeed \(q\geq Q^*_{r,s}\) implies \(q\geq2(r+s)-3\), and hence
\(r\leq(q+3)/2-s\leq |S|-\lfloor(q-1)/2\rfloor\).  The theorem therefore
says that these omitted curve points are all the weight-\(r\) directions.

Deleting columns cannot lower the weight of a tangent or conjugate-secant
direction.  Its full-support weight is \(r-1\), and it is not one of the
weight-\(r\) directions, so its \(S\)-weight remains \(r-1\).  Let \(f\) be an
interior point of a rational secant with at least one endpoint in \(A\).  A
representation of \(f\) by at most \(r-2\) retained columns, combined with
the secant representation, would give a nontrivial dependence among at most
\(r\) distinct normal-rational-curve points.  This contradicts the arc
property.  Hence \(d_S(f)\geq r-1\), and the deep-shell classification gives
equality.  Retained curve points have weight one, and interiors of secants
whose two endpoints are retained have weight two.  The escape theorem deals
with every point outside \(\cP_r\), so the list is exhaustive.

Tangent and conjugate-secant directions contribute
\(q(q+1)^2/2\).  The number of rational secants incident with \(A\) is
\[
 \binom{q+1}{2}-\binom{q+1-s}{2}=\frac{s(2q+1-s)}2,
\]
and each has \(q-1\) interior rational points.  This gives \(N_{r-1}\).
Syndrome vectors parametrize cosets, and every nonzero projective direction
has \(q-1\) representatives, giving the literal coset counts.
\end{proof}

For the full affine support \(A=\{\infty\}\), the theorem reads
\[
 N_r=1,\qquad N_{r-1}=\frac{q(q^2+4q-1)}2.
\]
The first equality is classical in this high-rate range; the second records
the complete next shell.

\subsection{MDS and NMDS appended columns}

For a nonzero syndrome column \(f\), put
\[
 \widehat C_f=\ker[H_S\mid f].
\]

\begin{corollary}[One-column extensions]
\label{cor:one-column-extensions}%
\coverage{absent}%
For every \(f\ne0\),
\begin{equation}\label{eq:extension-distance}
 d(\widehat C_f)=d_S(f)+1.
\end{equation}
Under the large-characteristic hypotheses of Theorem~\ref{thm:main}, the
omitted curve points give exactly the MDS extensions; the tangent,
conjugate-secant, and \(A\)-incident split-secant points give exactly the
NMDS extensions; every other appended column has Singleton defect at least
two.  For \(s=0\), the MDS class is empty.
\end{corollary}

\begin{proof}
The old columns of \(H_S\) have no dependent set of at most \(r\) columns.
A smallest dependence in \([H_S\mid f]\) either uses only old columns and
has size at least \(r+1\), or contains \(f\); after deleting \(f\), the latter
is a smallest old-column span of \(f\).  This proves
\eqref{eq:extension-distance}.

If \(d_S(f)=r-1\), the primal extension has distance \(r\), one below the
Singleton bound.  Choose \(r-1\) old columns spanning \(f\).  A nonzero
functional vanishing on their hyperplane gives a dual word of weight
\(|S|+1-r\), also one below the dual Singleton bound; the GRS distance of the
old part gives the reverse inequality.  Hence the extension is NMDS.
Conversely, an NMDS extension has primal distance \(r\), so
\eqref{eq:extension-distance} forces \(d_S(f)=r-1\).  The theorem now gives
the stated classes.
\end{proof}

\subsection{Family-aggregate NMDS enumerators}

Put \(m=r-1\), \(N=|S|=q+1-s\), and define
\[
 \mu_S(f)=\#\{T\subseteq S:|T|=m,\quad
                    f\in\langle\nu_m(T)\rangle\}.
\]
For a weight-\((r-1)\) projective syndrome direction \([f]\), every such
support gives one minimal dependence, unique up to a nonzero scalar.  Thus
\begin{equation}\label{eq:nmds-minimum-coefficient}
 A_r(\widehat C_f)=(q-1)\mu_S(f).
\end{equation}

Let \(\mathrm{Tan}\), \(\mathrm{Conj}\), and \(\mathrm{Split}_A\) denote
the sets of projective tangent, conjugate-secant, and deleted-point-incident
split-secant directions in this shell.  Sums use one representative of each
projective direction.

\begin{theorem}[Family-aggregate minimum supports]
\label{thm:family-aggregate-nmds}%
\coverage{absent}%
Under the large-characteristic hypotheses of Theorem~\ref{thm:main},
\begin{align}
 \sum_{f\in\mathrm{Tan}}\mu_S(f)
   &=(q+1-m)\binom Nm,\label{eq:aggregate-tangent}\\
 \sum_{f\in\mathrm{Conj}}\mu_S(f)
   &=\frac{q(q-1)}2\binom Nm,\label{eq:aggregate-conjugate}\\
 \sum_{f\in\mathrm{Split}_A}\mu_S(f)
   &=\left[\binom s2+s(N-m)\right]\binom Nm.
   \label{eq:aggregate-split}
\end{align}
These identities determine the complete family-aggregate and family-average
NMDS weight enumerators.  They do not assert that individual extensions in
one family have equal weight distributions.
\end{theorem}

\begin{proof}
For \(T\in\binom Sm\), let \(H_T\) be the hyperplane spanned by its normal-
rational-curve points.  On the tangent at \(a\), \(H_T\) meets the tangent at
the curve point if \(a\in T\), and at one off-curve point otherwise.  Thus
each \(T\) contributes \(q+1-m\) tangent incidences.  Every conjugate secant
meets \(H_T\) in exactly one rational point, giving \(q(q-1)/2\) incidences.
A secant through two omitted points contributes for every \(T\); a secant
through one omitted and one retained point contributes exactly when the
retained endpoint is outside \(T\).  These give respectively
\(\binom s2\) and \(s(N-m)\) incidences per \(T\).  Summing over all
\(\binom Nm\) sets \(T\) proves the three formulas.

If \(\mathcal A_j^{\mathcal F}=\sum_{[f]\in\mathcal F}
A_j(\widehat C_f)\), then \eqref{eq:nmds-minimum-coefficient} gives
\begin{align*}
 \mathcal A_r^{\mathrm{Tan}}
   &=(q-1)(q+1-m)\binom Nm,\\
 \mathcal A_r^{\mathrm{Conj}}
   &=(q-1)\frac{q(q-1)}2\binom Nm,\\
 \mathcal A_r^{\mathrm{Split}_A}
   &=(q-1)\left[\binom s2+s(N-m)\right]\binom Nm.
\end{align*}

Let \(L=N+1\) and \(K=L-r\).  The standard NMDS recurrence
\cite[Thm.~10]{MeneghettiPellegriniSala2020} expresses every coefficient as
\begin{align*}
 A_{r+i}(\widehat C_f)
={}&\binom{L}{K-i}
 \sum_{j=0}^{i-1}(-1)^j\binom{r+i}{j}(q^{i-j}-1)\\
 &+(-1)^i\binom Ki A_r(\widehat C_f),\qquad 1\leq i\leq K.
\end{align*}
Summing this identity over a family and applying
\eqref{eq:nmds-minimum-coefficient} proves the final assertion.
\end{proof}
